% Post-CP27 candidate; no promotion before complete independent audits.
\subsection{A candidate uniform nodal norm-eleven self-isogeny}
\label{subsec:uniform-eleven-theta-candidate}
\status{Candidate; complete proof and exact certificate pending independent audit.}

Work over an algebraically closed field of characteristic two.  The
normal form, Abel--Jacobi base point $W$, and quaternion automorphisms
are those of Lemma~\ref{lem:uniform-seven-interpolation-candidate}.
Write
\[
 C_z:\ y^2+y=ax^5+bx^3,\qquad a=z^2,\quad b=z^2+z+1.
\]
Choose a parameter $\xi\notin\mathbf F_4$, and set
\[
 r=\xi^2+\xi,\quad D=r^2,\quad t=\xi^2,\quad
 z=(r^2+1)/r^3.
 \tag{E11.1}
\]
Thus $r\ne0,1$, $t^2+t=D$, and all quaternion parameter denominators
are nonzero. Choose $\beta,w$ with
$\beta^2+\beta=z+1$ and $w^2+w=at^5+bt^3$, and put
\[
 \begin{split}
 G(x)&=zx^2+(z+1)x+\beta,\\
 H(x)&=z(t+r)x^2+(zt^2+t+r^{-2})x+w.
 \end{split}
\]
The automorphisms $g(x,y)=(x+1,y+G(x))$ and
$h(x,y)=(x+t,y+H(x))$ have $g^2=h^2=-1$ and $gh=-hg$ on the
Jacobian. Consequently
\[
 \psi=3+g+h,\qquad \psi^\dagger\psi=[11],\qquad \deg\psi=121
 \tag{E11.2}
\]
is an etale polarized self-isogeny. The integer $3$ here is an
endomorphism coefficient, not a function-field scalar.

\begin{lemma}[Exact interpolation and the finite theta kernel]
\label{lem:uniform-eleven-interpolation-candidate}
\textup{[Candidate.]}
Put
\[
 \gamma=\frac{\xi^{10}+\xi^7+\xi^5+\xi^3+1}{(\xi+1)^4},
 \qquad \alpha=\frac{t}{D}=\frac1{(\xi+1)^2},\quad
 \beta_0=\alpha+1.
\]
Suppose $\gamma\ne0$. For $P=(p,y_0)$ finite, define
\[
\begin{gathered}
 d_1=ap^4+bp^2,\quad d_2=d_1^2+bp,\quad
 Y=H(p)+t^2G(p)+Dd_1,\\
 Z=Y+tD\bigl(G(p)+d_1+d_2\bigr),\quad C=tY,\\
 A_1=Cd_1,\quad A_2=Zd_1+Cd_2,\quad
 A_3=Z(G(p)+d_1)+C(G(p)+d_1+d_2).
\end{gathered}
\]
Then, with $X=x+p$,
\[
 f_P=(ZX+C)(y+y_0)+A_1X+A_2X^2+A_3X^3
\]
has divisor $3P+gP+hP+E_P-7W$, where $E_P$ is an effective
nonspecial divisor of degree two. Moreover
$\ker\psi\cap\Theta=\{W\}$, and
\[
 Z=\alpha Y^2+\beta_0Y+\gamma.                   \tag{E11.3}
\]
Its norm is
\[
 f_P\iota(f_P)=X^3(X+1)(X+t)
       \bigl(aZ^2X^2+R_bX+R_c\bigr),            \tag{E11.4}
\]
where
\[
\begin{split}
 R_b&=A_3^2+a(p+t+1)Z^2,\\
 R_c&=\bigl(b+a((t+1)p+(t+1)^2+t)\bigr)Z^2
            +aC^2+(t+1)A_3^2.
\end{split}
\]
\end{lemma}
\begin{proof}
The quaternion assertions are the previously proved parameter identities;
Rosati sends $g,h$ to their negatives, so (E11.2) follows by multiplication.
For interpolation use the basis $1,x,x^2,x^3,y,xy$ of $H^0(7W)$.
An arbitrary section with order at least three at $P$ is
\[
 (lX+c)(y+y_0)+cd_1X+(ld_1+cd_2)X^2+kX^3.
\]
Its value at $gP$ determines
$k=(l+c)(G(p)+d_1)+cd_2$. The remaining value at $hP$ is exactly
\[
                 tYl+Zc=0.                    \tag{E11.5}
\]
These row operations are reversible. Identity (E11.3), obtained by
substitution using the three Artin--Schreier parameter relations, shows
that $Y,Z$ cannot both vanish. Since $t\ne0$, the solution space is
one-dimensional, spanned by $(l,c)=(Z,tY)$ and the displayed $f_P$.
Thus $h^0(7W-3P-gP-hP)=1$. Its effective residual divisor is nonspecial.
If $\psi(P-W)=0$, that space instead has dimension
$h^0(2W)=2$, a contradiction. Expansion gives (E11.4).
All these row and coefficient identities are checked universally in the
retained interpolation program, including the arbitrary-section row
(E11.5), not merely the one proposed section.
\end{proof}

\begin{lemma}[A five-dimensional algebra controlling every finite fibre]
\label{lem:uniform-eleven-quintic-fibres-candidate}
\textup{[Candidate.]}
Use the preceding parameters. Set
\[
\begin{gathered}
 a_4=aD,\quad a_2=(r^6+r^2+1)/r^4,\quad a_1=aD^2,\quad
 Y_0=w+t^2\beta,\\
 Q=at^5+bt^3+t^2(z+1),\quad
 A(L)=a_4L^3+a_2L+a_1,\quad
 L_0=aD^2(L+t+1)+Q,\\
 Z(Y)=\alpha Y^2+\beta_0Y+\gamma,\quad
 C_L(Y)=(Y+L_0)Z(Y)^2+Y^4.
\end{gathered}                                                  \tag{E11.6}
\]
The polynomial $C_L$ has degree five and is separable for every finite
$L$. In its quotient algebra, put
\[
\begin{split}
 u&=\frac{t^2}{\gamma^2}
   \left(\alpha^2Y^3+(L_0\alpha^2+1)Y^2
                    +\beta_0^2Y+L_0\beta_0^2\right),\\
 B&=Y+a_4(L^2u+u^2)+a_2u,\qquad
 P=B^2+LA(L)B+A(L)^2u.
\end{split}                                                     \tag{E11.7}
\]
Let
\[
 \chi_P(T)=T^5+c_4T^4+c_3T^3+c_2T^2+c_1T+c_0
\]
be the characteristic polynomial of multiplication by $P$, and let
$\chi_B$ be the analogous polynomial for $B$. A fibre with at least
three distinct points implies at least one of the following finite-$L$
conditions:
\[
\begin{array}{ll}
 c_3\ne0:&
 F_1=c_2^4+c_3^3c_2^2+c_1c_3^4=0,\\
 &F_2=c_2^5+c_4c_2^4c_3+c_1c_2c_3^4+c_0c_3^5=0;\\[2pt]
 c_3=0:&c_2=0;\\[2pt]
 A(L)=0:&\operatorname{Res}_T(\chi_B,\partial_T\chi_B)=0.
\end{array}                                                     \tag{E11.8}
\]
\end{lemma}
\begin{proof}
First remove the exceptional coordinates. Infinity $W$ has no other
preimage by the preceding lemma. Two finite points with equal
$x$-coordinate are hyperelliptic conjugates. Equality of their images
would make that image geometric two-torsion; in two-rank zero it would
be zero, already excluded. Thus at most one point above each finite
$x$-coordinate occurs in any fibre.

At $Z=0$, (E11.3) has two distinct nonzero roots in $Y$, since
$\beta_0=t^2/D\ne0$ and $\gamma\ne0$. The interpolation identity
$DA_3=Z(tG(p)+H(p))+Y^2$ gives $A_3=Y^2/D\ne0$.
The function has pole order six, so $E_P=W+R$ with
\[
 x(R)=p+t+1+at^2D^2/Y^2.
\]
For each of the two possible $Y$'s a fixed $R$ determines at most one
$p$. Hence these fibres have at most two points. They cannot coincide
with a fibre at $Z\ne0$, whose unique nonspecial residual divisor is
finite of degree two. This argument also allows a doubled finite
residual divisor.

For $Z\ne0$, write the residual monic polynomial in the global
$x$-coordinate as $x^2+Lx+M$. Directly from (E11.4), and using
\[
 Y=a_4p^4+a_2p^2+a_1p+Y_0,\quad
 (tG(p)+H(p))^2=aD^2p+Y+Q,
\]
one obtains
\[
 C_L(Y)=0,\qquad
 U:=M+(t+1)L+t+b/a=p^2+Lp+t^2Y^2/Z(Y)^2.       \tag{E11.9}
\]
The derivative of $C_L$ with respect to $Y$ is $Z(Y)^2$.
A common zero of $C_L,Z$ would have $Y=0$, contradicting
$Z(0)=\gamma$. Thus its five roots are distinct for every finite $L$.
The leading coefficient $\alpha^2$ is a nonzero parameter constant.
The polynomial $u$ in (E11.7) satisfies
$uZ(Y)^2=t^2Y^2$ in this algebra. This identity is checked after
polynomial division by the monic $\alpha^{-2}C_L$; it holds also at
$Y=0,L_0=0$, so no inversion of $Y$ or $L_0$ is assumed.

Put $J=Y_0+a_4(L^2U+U^2)+a_2U$.
Reducing the quartic expression for $Y$ by $p^2+Lp=U+u$ gives
\[
                  A(L)p=B+J.                  \tag{E11.10}
\]
If $A(L)\ne0$, a fixed pair $(U,Y)$ determines $p$ uniquely, and
\[
              P(Y)=J^2+LA(L)J+A(L)^2U.
\]
Three points in a fibre therefore give three distinct roots of $C_L$
with one common $P$-value. That value is a triple root of $\chi_P$.
Its first ordinary derivative is $T^4+c_3T^2+c_1$ and its second
Hasse derivative is $c_3T+c_2$. If $c_3\ne0$, substitution of
$T=c_2/c_3$ gives $F_1,F_2$; if $c_3=0$, it requires $c_2=0$.
Omitting the additional equation actually needed for a triple root
gives a weaker necessary condition, hence a larger bad-parameter locus
in the certificate.

If $A(L)=0$, (E11.10) requires $B(Y)=J$. When the five $B$-values
are distinct, this determines at most one $Y$. Equation (E11.9)
then has at most two solutions for $p$, including when $L=0$.
Thus three points force the last resultant in (E11.8) to vanish.
No claim about separability of that quadratic is needed.
\end{proof}

\begin{lemma}[An exact parameter-avoidance certificate]
\label{lem:uniform-eleven-good-parameters-candidate}
\textup{[Candidate; exact universal certificate supplied.]}
For every normalized modulus $\mu$, there is a parameter $\xi$ as
in (E11.1) with $\gamma\ne0$ for which none of (E11.8) occurs.
\end{lemma}
\begin{proof}
Here are the complete finite algebraic inputs and the uniform inference.
All polynomials are over $\mathbf F_2$. Put
\[
 H_4(V)=V^4+V,\quad D_{12}(V)=V^{12}+V^9+V^6+V^3+1,
 \quad G_{10}(V)=V^{10}+V^7+V^5+V^3+1.
\]
The characteristic coefficients are computed in
$\mathbf F_2[V,V^{-1},(V+1)^{-1},G_{10}^{-1}][L]$ from the monic
quintic and the two multiplication matrices of (E11.6)--(E11.7).
Each coefficient is retained as a list of binary polynomials in $V$
indexed by the power of $L$, together with its three denominator
exponents. Thus the inputs can be reconstructed without numerical
sampling or an implicit choice of roots.

Clear denominators coefficientwise in $F_1,F_2,c_3,c_2,A$ to their
retained numerators $\widetilde F_1,\widetilde F_2,
\widetilde c_3,\widetilde c_2,\widetilde A$.
For $\chi_B$, use the common denominator
$V^{10}(V+1)^{18}G_{10}^8$, and call the resulting degree-five
polynomial $\widetilde\chi_B$. Define
\[
\begin{split}
 R_1&=\operatorname{Res}_L(\widetilde F_1,\widetilde F_2),\\
 R_2&=\operatorname{Res}_L(\widetilde c_3,\widetilde c_2),\\
 R_3&=\operatorname{Res}_L\left(\widetilde A,
       \operatorname{Res}_T(\widetilde\chi_B,
                           \partial_T\widetilde\chi_B)\right).
\end{split}                                                     \tag{E11.11}
\]
The exact computations give nonzero polynomials of degrees
$121104,3944,3540$, respectively. Their full binary coefficients and
complete factor lists are retained. Let $B_*(V)$ be the product, once
each, of the distinct monic factors in those lists, together with
$G_{10}$ if needed, after deleting factors of $H_4$.
It has degree $3354$. Factor-product identities verify all three
resultants. Irreducibility of the listed factors is not required for
the zero-set implication being used.

Write $B_* = B_0+B_1+B_2$, where $B_j$ collects terms whose exponent
is $j$ modulo three. The second exact identity is
\[
 \gcd\left(B_*(V),\ B_*(V+1),\ V^{3354}B_*(1/V),
                  B_0+B_2,\ B_1+B_2\right)=1.                  \tag{E11.12}
\]
A retained univariate Bezout identity provides an independently
checkable form of (E11.12). These are universal polynomial operations,
not a census of finitely many surfaces.

The presentation modulus is
\[
 \mu^5=\frac{(z^2+z+1)^5}{z^6}
       =\frac{D_{12}(\xi)^5}{H_4(\xi)^{12}}.                  \tag{E11.13}
\]
For any given $\mu$, the numerator equation of (E11.13) is monic of
degree sixty and has a root. No such root belongs to $\mathbf F_4$,
since $D_{12}H_4=V^{16}+V$ and $D_{12}$ is nonzero at every root
of $H_4$. Fix one root $\xi$, and choose
$\omega\in\mathbf F_4$ with $\omega^2+\omega+1=0$.
All five parameters
\[
             \xi,\quad \xi+1,\quad \xi^{-1},
                    \quad \omega\xi,\quad \omega^2\xi
\]
remain outside $\mathbf F_4$ and have the same value of (E11.13).
This follows directly from translation invariance of $D_{12},H_4$,
reciprocity of $D_{12}$, and its exponents being divisible by three.
No assertion about a generic orbit size is needed, including at
$\mu=0$.

If all five parameters were bad, $\xi$ would annihilate the first
three polynomials of (E11.12), and both $B_*(\omega\xi)$ and
$B_*(\omega^2\xi)$. The latter two equations are equivalent to
$(B_0+B_2)(\xi)=(B_1+B_2)(\xi)=0$, contradicting (E11.12).
Thus one parameter avoids $B_*$. All deleted parameter factors are
already excluded by $\xi\notin\mathbf F_4$, while $\gamma\ne0$
is guaranteed by the included $G_{10}$. Nonvanishing of the three
resultants (E11.11) excludes every condition (E11.8), including all
specializations where a leading coefficient in $L$ drops.
Finally the choice of fifth-root normalization in the previously
proved normal-form argument gives the prescribed $\mu$, not merely
its fifth power.
\end{proof}

\begin{theorem}[A uniform nodal eleven-self-isogeny]
\label{thm:uniform-eleven-nodal-self-isogeny-candidate}
\textup{[Candidate.]}
Every indecomposable two-rank-zero principally polarized abelian surface
in characteristic two has an etale polarized self-isogeny of multiplier
$11$ whose principal theta image has exactly $120$ ordinary nodes and
no other singularities.
\end{theorem}
\begin{proof}
Use the good presentation in the preceding lemma and (E11.2).
The fibre lemma excludes every three-point fibre, with all exceptional
coordinates treated explicitly. Lemma~\ref{lem:root-nonordinary-odd-theta-nodality}
therefore applies: an etale odd quotient of a two-rank-zero genus-two
Jacobian with fibres of size at most two has only transverse ordinary
nodes. Its node count is $|\ker\psi|-1=121-1=120$.
\end{proof}

\begin{corollary}[All odd degrees divisible by eleven in two-rank zero]
\label{cor:f0-odd-multiples-eleven-candidate}
\textup{[Candidate.]}
On every fixed two-rank-zero principally polarized surface over a perfect
field of characteristic two, every positive odd degree divisible by at
least one of $3,5,7,11$ has relative geometric and arithmetic image
$\operatorname{GL}_{9d^2-1}(\mathbf F_2)$. Every even degree already
has the proved full one-covector stabilizer. The first remaining odd
fixed-surface degree in two-rank zero is $13$.
\end{corollary}
\begin{proof}
Theorem~\ref{thm:good-theta-odd-norm-step-candidate} gives degree eleven
from degree one on every
indecomposable two-rank-zero surface. For $d=11n$, choose a nonproduct
incoming source of odd multiplier $n$ by
Theorem~\ref{thm:nonordinary-arbitrary-odd-neighbor-candidate}; it is
again two-rank zero and has the newly proved degree-eleven group.
Apply Theorem~\ref{thm:cubic-component-arbitrary-odd-norm-transfer-candidate}
with $e=11$ when $n>1$. Products and $n=1$ are the previous product
theorems and the base just established. The $3,5,7$ and even assertions
were already proved. Full geometric general linear image forces the
same arithmetic image. No isogeny is assumed to be defined over the
original perfect field.
\end{proof}

\begin{remark}[Verification and remaining scope]
\label{rem:uniform-eleven-verification-scope-candidate}
\textup{[Candidate; audit pending.]}
The exact universal interpolation is retained as
\path{code/monodromy/post_CP27_eleven_verified_interpolation.py};
the rational-quintic reconstruction is
\path{code/monodromy/post_CP27_eleven_rational_quintic.py}, using
\path{code/monodromy/post_CP27_binary_coefficient_ring.py}.
The full generic and exceptional resultants, the $A=0$ discriminant
resultant, and the five-parameter gcd are retained respectively in
\path{code/monodromy/post_CP27_eleven_flint_elimination.py},
\path{code/monodromy/post_CP27_eleven_A_zero_elimination.py}, and
\path{code/monodromy/post_CP27_eleven_bad_parameter_orbits.py}.
The independent factor-product and explicit univariate Bezout reconstruction
is \path{code/monodromy/post_CP27_eleven_certificate_reconstruction.py};
it also checks $uZ^2=t^2Y^2$ directly without inverting $Y$.
Its first complete author run passed at 23:25:07 UTC on 19 November 2026
in 55.395332292 seconds.
Their complete outputs are in the matching
\path{agents/monodromy/outputs/} files. The resultant discovery uses
\texttt{python-flint 0.9.0}; all coefficient inputs, outputs, factor
products, and the univariate Bezout data are retained for independent
reconstruction. Earlier bounded timeouts, a serialization failure, and
an exploratory finite-field inference that omitted the condition
$c_3\ne0$ are preserved as failed attempts. None was used in a
promoted theorem. This candidate makes no new two-rank-one claim.
\end{remark}
